% !TEX root = sga4.5-en.tex
% ENGLISH EDITION -- TRANSLATED FROM FRENCH (AI-assisted, needs human review).

\chapter{Duality}\label{V}



















This expos\'e contains a few theorems and compatibilities, all
concerning Poincar\'e duality. In \ref{V:1}, the theorem of
local biduality in dimension $1$ \cite[I.5.1]{sga5} and a few
related calculations. In \ref{V:2}, a very economical proof of
Poincar\'e duality for curves, communicated to me by M.\ Artin. In
\ref{V:3}, a compatibility linking two definitions of the
pairing giving rise to Poincar\'e duality for curves: via the
cup-product, or via self-duality of the Jacobian. In \ref{V:4},
finally, the proof of the compatibility of the title.

Throughout this expos\'e, schemes will be noetherian and separated, and $n$
is an integer invertible on all the schemes considered.










\section{Local biduality, in dimension \texorpdfstring{$1$}{1}}\label{V:1}




\subsection{}\label{V:1-1}

Let $S$ be a regular scheme purely of dimension $1$. We propose
to show that the complex reduced to $\dZ/n$ in degree $0$ is
dualizing, i.e.\ that for $\sK\in \ob\eD_c^b(S,\dZ/n)$ ($(-)_c$ for
constructible), if we put $D\sK=\rHom(\sK,\dZ/n)$, then $D\sK$ is again
constructible with bounded cohomology, and that the canonical morphism $\alpha$
from $\sK$ to $D D\sK$ is an isomorphism.

For $\sK$ in $\eD^-$, and $\sL$ arbitrary, we have
\[
  \hom(\sK\lotimes \sL,\dZ/n) = \hom(\sL,\Hom(\sK,\dZ/n)) \text{.}
\]
The morphism from $\sK$ to $DD\sK$ is defined assuming $D\sK$ in
$\eD^-$; if $\beta:D\sK\lotimes\sK\to \dZ/n$ is the canonical pairing, it
is defined by the pairing
$\sK\lotimes D\sK=D\sK\lotimes \sK \xrightarrow\beta \dZ/n$.




\subsection{}\label{V:1-2}

Let $f:X\to S$ be a separated morphism of finite type. Put
$\sK_X=\eR f^!\dZ/n$. For $\sK\in \ob\eD^-(X,\dZ/n)$, put
$D\sK = \rHom(\sK,\sK_X)$. Adjunction between $\eR f_!$ and $\eR f^!$
ensures that
\[\xymatrix{
  \eR f_\ast \sK\lotimes \eR f_\ast D\sK \ar[r]
    & \eR f_! (\sK\lotimes D\sK) \ar[r]
    & \eR f_! \eR f^! \dZ/n \ar[r]
    & \dZ/n \text{.}
}\]
The symmetry of this description shows that, for $f$ proper, the diagram
\begin{equation*}\tag{1.2.1}\label{V:eq:1-2-1}
\xymatrix{
  \eR f_\ast \sK \ar[r] \ar[d]
    & D D \eR f_\ast \sK \ar[d]^-\wr \\
  \eR f_\ast D D \sK \ar[r]^\sim
    & D\eR f_\ast D\sK
}
\end{equation*}
is commutative.

If $f$ is the inclusion of a closed point, we know that
$\eR f^!\dZ/n=\dZ/n(-1)[-2]$ -- i.e., up to twist and shift, $\dZ/n$.
On the spectrum of a field, this complex is dualizing (Pontrjagin duality
for $\dZ/n$-modules). By \eqref{V:eq:1-2-1}, we hence have
$\sK\iso D D\sK$ for $\sK$ of the form $\eR f_!\sL$ -- and hence when the
support of $\underline\h^\bullet(\sK)$ is finite.




\begin{theorem_}\label{V:1-3}
Let $j:U\hookrightarrow S$ be a dense open of $S$ and $\sF$ a constructible
locally constant sheaf of $\dZ/n$-modules on $U$. We have
$\eD j_\ast \sF = j_\ast D \sF$, i.e.\ $\Hom(j_\ast \sF,\dZ/n) = j_\ast\Hom(\sF,\dZ/n)$ and $\Ext^i(j_\ast\sF,\dZ/n) = 0$ for $i>0$.
\end{theorem_}

On $U$, $\Ext^i(j_\ast\sF,\dZ/n)=0$ for $i>0$, since $\sF$ is locally
constant (and $\dZ/n$ is an injective $\dZ/n$-module), whereas for $i=0$
it is the dual $\sF^\vee$ of $\sF$. One checks that
$\hom(j_\ast\sF,\dZ/n)=j_\ast\sF^\vee$, and it remains to check the vanishing
of the $\Ext^i$ ($i>0$) at the points of $S\setminus U$. The problem is local
at these points. This reduces us to assuming that $S$ is a strictly local
trait and that $U$ is reduced to its generic point $\eta$. Let
$I=\gal(\bar\eta/\eta)$. The sheaf $\sF$ identifies with the Galois module
$\sF_{\bar\eta}$, and the special fibre of $j_\ast \sF$ with
$\sF_{\bar\eta}^I$.

Let $i$ be the inclusion of the closed point $s$, and apply $D$ to the exact
sequences
\[\xymatrix{
  0 \ar[r]
    & j_! \sF \ar[r]
    & j_\ast \sF \ar[r]
    & i_\ast \sF_{\bar\eta}^I \ar[r]
    & 0 \\
  0 \ar[r]
    & j_! \sF_{\bar\eta}^I \ar[r] \ar[u]
    & \sF_{\bar\eta}^I \ar[r] \ar[u]
    & i_\ast \sF_{\bar\eta}^I \ar[r] \ar[u]
    & 0 \text{.}
}\]

We obtain a morphism of triangles
\[\xymatrix{
  i_\ast\left(\sF_{\bar\eta}^I\right)^\vee (-1)[-2] \ar[r] \ar@{=}[d]
    & D j_\ast\sF \ar[r] \ar[d]
    & \eR j_\ast \sF^\vee \ar[d] \\
  i_\ast\left(\sF_{\bar\eta}^I\right)^\vee(-1)[-2] \ar[r]
    & \left(\sF_{\bar\eta}^I\right)^\vee \ar[r]
    & \eR j_\ast \sF_{\bar\eta}^I \text{.}
}\]
The long exact sequence deduced from the first row gives the vanishing of the
$\Ext^i$ ($i>2$) and, taking the fibre at $s$, we find
\[\xymatrix@=0.7cm{
  0 \ar[r]
    & \left(\underline\h^1(D j_\ast\sF\right)_s \ar[r]
    & \h^1\left(I,\sF_{\bar\eta}^\vee\right) \ar[r]^-\partial \ar[d]
    & \left(\sF_{\bar\eta}^I\right)^\vee(-1) \ar[r] \ar@{=}[d]
    & \left(\underline\h^2 D j_\ast\sF\right)_s \ar[r]
    & 0 \\
  & 0 \ar[r]
    & \h^1\left(I,(\sF_{\bar\eta}^I)^\vee\right) \ar[r]^-\partial
    & \left(\sF_{\bar\eta}^I\right)^\vee(-1) \ar[r]
    & 0 \text{.}
}\]
Since
$\left(\sF_{\bar\eta}^I\right)^\vee = \left(\sF_{\bar\eta}^\vee\right)_I$,
putting $M=\sF_{\bar\eta}^\vee$, it finally remains to check that
\[\xymatrix{
  \h^1(I,M) \ar[r]^-\sim
    & \h^1(I,M_I) \text{.}
}\]
If $p$ is the residual characteristic exponent, $I$ is an extension of a
group isomorphic to $\widehat\dZ_{p'} = \varprojlim_{(m,p)=1} \dZ/m$ by a
$p$-group $P$. Since $P$ is prime to the order of $M$, we have
$\h^1(P,M)=0$ ($i>0$) and $(M_I)^P=(M^P)_I$. This allows us to replace $M$ by
$M^P$ and $I$ by $I/P$. Finally we have a functorial isomorphism
$\h^1(\widehat\dZ_{p'},M)\sim ($coinvariants of $\widehat\dZ_{p'}$ in $M)$,
whence the theorem.




\begin{theorem_}\label{V:1-4}
For $\sK\in \ob\eD_c^b(S,\dZ/n)$, we have $\sK\iso D D \sK$.
\end{theorem_}

By d\'evissage, we reduce to assuming that $\sK$ is reduced to a
constructible sheaf $\sF$ in degree $0$, and that $\sF$ is either with finite
support (\ref{V:1-2}), or of the form $j_\ast\sF_1$ as in \ref{V:1-3}. In
this second case, \ref{V:1-3} reduces us to local biduality for $\sF$
locally constant on $U$.










\section{Poincar\'e duality for curves, after M.\ Artin}\label{V:2}




\subsection{}\label{V:2-1}

Let $X$ be a smooth projective curve over $k$ algebraically closed. Put
$\sK_X=\dZ/n(1)[2]$ and, for $\sK\in\ob \eD_c^b(X,\dZ/n)$,
$D\sK=\rHom(\sK,\sK_X)$. For $M$ a $\dZ/n$-module, also put
$D M = \hom(M,\dZ/n)$; likewise for complexes of modules. The trace
morphism $\h^0(X,\sK_X) \to \dZ/n$, or $\eR\Gamma(X,\sK_X)\to \dZ/n$, defines a
pairing $\eR\Gamma(X,\sK)\lotimes\eR\Gamma(X,D\sK) \to \dZ/n$. Poincar\'e
duality between cohomology and cohomology with proper supports of an open
$j:U\hookrightarrow X$ of $X$ says that, for $\sK=j_!\dZ/n$, this pairing
identifies each factor with the dual of the other.

I explain below a proof, communicated to me by
M.\ Artin, that for $\sK\in\ob\eD_c^b(X,\dZ/n)$, this pairing is
always perfect, i.e.\ defines an isomorphism
\begin{equation*}\tag{2.1.1}\label{V:eq:2-1-1}
\xymatrix{
  \eR\Gamma(X,\sK) \ar[r]^-\sim
    & D \eR\Gamma(X,D\sK) \text{.}
}
\end{equation*}

For every constructible sheaf $\sF$, put
\[
  '\h^i(X,\sF) = \text{$\dZ/n$-dual of $\h^{-i}(X,D\sF)$.}
\]
That \eqref{V:eq:2-1-1} be an isomorphism is equivalent to the




\begin{theorem_}\label{V:2-2}
For $\sF$ a constructible sheaf of $\dZ/n$-modules, we have
\begin{equation*}\tag{2.2.1}\label{V:eq:2-2-1}
\xymatrix{
  \h^i(X,\sF) \ar[r]^-\sim
    & '\h^i(X,\sF) \text{.}
}
\end{equation*}
\end{theorem_}




\begin{lemma_}\label{V:2-3}
Let $f:X\to Y$ be a generically \'etale morphism between smooth curves over
$k$. We have $\sK_X=f^\ast\sK_Y=\eR f^!\sK_Y$, with $\tr_f$ as
adjunction arrow $\eR f_!\eR f^!\sK_Y\to \sK_Y$.
\end{lemma_}

% NOTE: \rhom or \rHom ? originally \rhom
We reduce to checking that for $f$ finite, we have
$f_\ast \eR f^!\sK_Y = \rHom(f_\ast\dZ/n,\sK_Y)$ is $f_\ast\dZ/n$, with $\tr_f$
as adjunction arrow. The vanishing of the $\Ext^i$ ($i>0$) follows from
\ref{V:1-3}, and the assertion follows.

One can say, more explicitly, that $\eR f^!$ is the derived functor of the
functor $f^!:f^!\sF(U)=\hom(f_!\dZ_U,\sF)$ and \ref{V:1-3} implies the
vanishing of the $\eR^i f^!\dZ/n$ for $i>0$.




\begin{corollary_}\label{V:2-4}
Let $f:X'\to X$ be a generically \'etale morphism of smooth
projective curves over $k$. We have $'\h^i(X,f_\ast \sF)='\h^i(X',\sF)$; this
isomorphism is functorial and the diagram
\[\xymatrix{
  \h^i(X,f_\ast\sF) \ar[r] \ar@{=}[d]
    & '\h^i(X,f_\ast\sF) \ar@{=}[d] \\
  \h^i(X',\sF) \ar[r]
    & '\h^i(X',\sF)
}\]
is commutative.
\end{corollary_}

The isomorphism is given by \ref{V:1-2}: $D f_\ast\sF=f_\ast D\sF$.




\subsection{}\label{V:2-5}

Let us prove \ref{V:2-2}. If $\sF$ is reduced to $\dZ/n$ at a point, extended
by $0$, we have $\h^\bullet(D\sF) = \ext^\bullet(\sF,\sK_X)=\dZ/n$ in degree
$0$, and in the pairing $\h^\bullet(\sF)\otimes\h^\bullet(D\sF) \to \dZ/n$,
we have $1\otimes 1\mapsto 1$ (\hyperref[IV]{Cycle}, \ref{IV:2-1-5}): the duality
is perfect. The case where $\sF$ has finite support is treated likewise.

For arbitrary constructible $\sF$, $\sA$ the subsheaf of its sections with
finite support, and $\sG=\sF/\sA$, the exact sequence
\begin{equation*}\label{V:eq:2-5-1}\tag{2.5.1}
\xymatrix{
  0 \ar[r]
    & \sA \ar[r]
    & \sF \ar[r]
    & \sG \ar[r]
    & 0
}
\end{equation*}
gives a diagram
\begin{equation*}\tag{2.5.2}\label{V:eq:2-5-2}
\xymatrix@=.5cm{
  & & 0 \ar[r]
    & \h^0(X,\sA) \ar[r] \ar[d]^-\wr
    & \h^0(X,\sF) \ar[r] \ar[d]
    & \h^0(X,\sG) \ar[r] \ar[d]
    & 0 \\
  0 \ar[r]
  & '\h^{-1}(X,\sF) \ar[r]
  & '\h^{-1}(X,\sF) \ar[r]
  & '\h^0(X,\sA) \ar[r]
  & '\h^0(X,\sF) \ar[r]
  & '\h^0(X,\sF) \ar[r]
  & 0
}
\end{equation*}
and compatible isomorphisms $\h^i(X,\sF)\iso \h^i(X,\sG)$,
$'\h^i(X,\sF) \iso '\h^i(X,\sG)$ ($i\ne 0,-1$). If $j:U\to X$ is an open on
which $\sF$ is locally constant, we have an exact sequence
\begin{equation*}\tag{2.5.3}\label{V:eq:2-5-3}
\xymatrix{
  0 \ar[r]
    & \sG \ar[r]
    & j_\ast j^\ast \sF \ar[r]
    & \sB \ar[r]
    & 0
}
\end{equation*}
with $\sB$ of finite support. Since
$'\h^i(X,j_\ast j^\ast\sF) = D \h^{2-i} (X,j_\ast(j^\ast\sF)^\vee(1))$
by \ref{V:1-3}, this $'\h^i=0$ for $i<0$ and \eqref{V:eq:2-5-2} and the
long exact sequence defined by \eqref{V:eq:2-5-3} show that $'\h^i$ always
vanishes for $i<0$.

If $\sF=\dZ/n$, the theorem holds for $i=0$ and $2$. This expresses that,
for $X$ connected, the trace morphism: $\h^2(X,\dmu_n) \iso \dZ/n$ is an
isomorphism. By \ref{V:2-4}, the theorem remains true for $i=0$ and
$\sF=f_\ast\dZ/n$.

For $\sF$ again arbitrary and $\sG$ as above, there exists
$f:X'\to X$ as in \ref{V:2-4} such that $\sG$ embeds into $f_\ast\dZ/n$:
\[\xymatrix{
  0 \ar[r]
    & \sG \ar[r]
    & f_\ast\dZ/n \ar[r]
    & \sQ \ar[r]
    & 0 \text{.}
}\]
Replacing $X'$ by an $X''/X'$, we may assume that
$\h^i(X,\sF) \to \h^i(X,f_\ast\dZ/n) = \h^i(X',\dZ/n)$ is zero for $i>0$.
Consider then
\small
\[\xymatrix@=.3cm{
  0 \ar[r]
    & \h^0(X,\sG) \ar[r] \ar[d]
    & \h^0(S',\dZ/n) \ar[r] \ar[d]^-\wr
    & \h^0(X,\sQ) \ar[r] \ar[d]
    & \h^1(X,\sG) \ar[r]^-0 \ar[d]
    & \h^1(X',\dZ/n) \ar[r] \ar[d]
    & \h^1(X,\sQ) \ar[r] \ar[d]
    & \cdots \\
  0 \ar[r]
    & '\h^0(X,\sG) \ar[r]
    & '\h^0(X',\dZ/n) \ar[r]
    & '\h^0(X,\sQ) \ar[r]
    & '\h^1(X,\sG) \ar[r]
    & '\h^1(X',\dZ/n) \ar[r]
    & '\h^1(X,\sQ) \ar[r]
    & \cdots
}\]
\normalsize

The isomorphism in $\h^0(X')$ gives $\h^0(X,\sG)\hookrightarrow '\h^0(X,\sG)$,
and the same injectivity for $\sF$ \eqref{V:eq:2-5-2}. Applied to $\sQ$,
this gives $\h^1(X,\sG)\hookrightarrow '\h^1(X,\sG)$. One applies this to
$\h^1(X',\dZ/n)$. This group having the same order as
$'\h^1(X',\dZ/n)=D(\h^1(X',\dZ/n)(1))$, one finds an isomorphism, and
$\h^1(X,\sG) \iso '\h^1(X,\sG)$. Likewise for $\sF$. For $i=2$, we already
know that $\h^2(X',\dZ/n) \iso '\h^2(X',\dZ/n)$, and we proceed likewise,
then stop.










\section{Poincar\'e duality for curves}\label{V:3}

In this section, we prove the compatibility (\hyperref[I]{Arcata}
\ref{I:eq:6-2-3-3}).




\subsection{}\label{V:3-1}

The notations are those of (\hyperref[I]{Arcata} \ref{I:6-2-3}): $\bar X$ is
a smooth projective connected curve over an algebraically closed field $k$,
$D$ is a reduced divisor on $\bar X$, $0$ a point of
$X=\bar X\setminus D$
\[\xymatrix{
  X \ar@{^{(}->}[r]^-j
    & \bar X
    & \ar[l]_-i D \text{,}
}\]
and $\pic_D(\bar X)=\h^1(\bar X,_D\dG_m)$ where $_D\dG_m$ is the sheaf of
sections of $\dG_m$ congruent to $1\bmod D$. Recall that
$\h^1(X,\dmu_n) = \pic_D^0(\bar X)_n$, and that $x\to \sO(x)$ defines a
canonical map $f:X\to \pic_D(\bar X)$. Put
$f_0(x)=f(x)-f(0):X\to \pic_D^0(\bar X)$.

Denote again by $j$ the inclusion of $X\times X$ in $X\times \bar X$. The
diagonal $\Delta$ of $X\times X$ is closed in $X\times \bar X$; it
defines a class in
$\h_\Delta^2(X\times X,\dmu_n)=\h_\Delta^2(X\times \bar X,j_!\dmu_n)$. Denote
by $c$ its image in $\h^2(X\times \bar X,j_!\dmu_n)$.

The K\"unneth formula ensures that
\[
  \h^\bullet(X\times \bar X,j_!\dmu_n) = \h^\bullet(X,\dZ/n)\otimes \h^\bullet(\bar X,j_!\dmu_n) = \h^\bullet(X,\dZ/n)\otimes \h_c^\bullet(X,\dmu_n) \text{.}
\]
We propose to compute the $(1,1)$-component of $c$,
\[
  c^{1,1} \in \h^1(X,\dZ/n)\otimes \h_c^1(X,\dmu_n) = \h^1(X,\h_c^1(X,\dmu_n)) = \h^1\left(X,\pic_D^0(\bar X)_n\right) \text{.}
\]

The exact sequence
\[\xymatrix{
  0 \ar[r]
    & \pic_D^0(\bar X)_n \ar[r]
    & \pic_D^0(\bar X) \ar[r]^n
    & \pic_D^0(\bar X) \ar[r]
    & 0
}\]
makes $\pic_D^0(\bar X)$ a $\pic_D^0(\bar X)_n$-torsor over
$\pic_D^0(\bar X)$. Denote by $u$ the class, in
$\h^1(X,\h_c^1(X,\dmu_n))=\h^1(X,\pic_D^0(\bar X)_n)$, of its inverse
image under $f_0$.




\begin{proposition_}\label{V:3-2}
With the above notations, $c^{1,1}=-u$.
\end{proposition_}

Let $e$ be the difference of the images in $\h^2(X\times \bar X,j_!\dmu_n)$ of the
cohomology classes of $\Delta$ and $X\times\{0\}$ being of K\"unneth
type $(0,2)$, we have $c^{1,1}=e^{1,1}$. Let the Leray spectral sequence
for $\pr_1:X\times \bar X\to X$ and the sheaf $j_!\dmu_n$ be
\[
  E_2^{pq} = \h^p(X,\eR^q {\pr_1}_\ast j_!\dmu_n) = \h^p(X,\h_c^q(X,\dmu_n)) \Rightarrow \h^{p+q}(X\times \bar X,j_!\dmu_n) \text{.}
\]
It degenerates: it is the tensor product of $\h^\bullet(X)$ with the
(trivial) Leray spectral sequence for $\bar X\to \spec(k)$ and the sheaf
$j_!\dmu_n$. The divisor $\Delta-X\times\{0\}$ is fibre by fibre of degree
$0$, so that the image of $e$ in $E^{02}$ is zero. We can hence speak
of its image $\bar e$ in $E^{1,1}$, and we must show that
$\bar e = -u$.

Let $\sG$ on $X\times\bar X$ be the subsheaf of $\dG_m$ consisting of local
sections whose restriction to the subscheme $X\times D$ is $1$. We again have
an exact sequence
\begin{equation*}\tag{3.2.1}\label{V:eq:3-2-1}
\xymatrix{
  0 \ar[r]
    & j_!\dmu_n \ar[r]
    & \sG \ar[r]
    & \sG \ar[r]
    & 0
}
\end{equation*}
and $e$ is the image under $\partial$ of the class $e_1\in \h^1(X\times\bar X,\sG)$
of the invertible sheaf $\sO(\Delta-X\times\{0\})$ trivialized by $1$ on
$X\times D$.

The direct image under $\eR\pr_{1\ast}$ of the distinguished triangle defined by
\eqref{V:eq:2-2-1} is a distinguished triangle
\begin{equation*}\tag{3.2.2}\label{V:eq:3-2-2}
\xymatrix{
  & \ar[r]^-\partial
    & \eR \pr_{1\ast} j_!\dmu_n \ar[r]
    & \eR \pr_{2\ast}\sG \ar[r]
    & \eR \pr_{2\ast}\sG \ar[r]^-\partial
    &
}
\end{equation*}
and the diagram
\[\xymatrix{
  \h^1(X\times \bar X,\sG) \ar[r]^-\partial \ar@{=}[d]^-{(1)}
    & \h^2(X\times \bar X,j_!\dmu_n) \ar@{=}[d] \\
  \h^1(X,\eR\pr_{1\ast} \sG) \ar[r]^-\partial
    & \h^2(X,\eR\pr_{1\ast} j_!\dmu_n)
}\]
is commutative: our problem becomes that of identifying the image in
$E^{11}=\h^1(X,\h_c^1(X,\dmu_n))$ (the image in $E^{02}$ being zero) of
the image under
$\partial_{\text{\eqref{V:eq:3-2-2}}}:\h^1(X,\eR\pr_{1\ast}\sG) \to \h^2(X,\eR\pr_{1\ast} j_!\dmu_n)$
of the class, still denoted $e_1$, corresponding to $e_1$ via
(1).

The sheaf $\eR^1\pr_{1\ast}\sG$ is defined by the group scheme over
$X$ inverse image of the algebraic group $\pic_D(\bar X)$ over
$\spec(k)$, and the image of $e_1$ in
$\h^0(X,\eR^1\pr_\ast\sG) = \hom(X,\pic_D(\bar X))$ is $f_0$.

Represent the triangle \eqref{V:eq:3-2-2} as defined by a short
exact sequence of complexes of sheaves.
\[\xymatrix{
  0 \ar[r]
    & \sA \ar[r]
    & \sB \ar[r]
    & \sC \ar[r]
    & 0 \text{.}
}\]
Let $\sB'$ be the subsheaf of
\[
  \tau_{\leqslant 1} \sB : \cdots \sB^0 \to \ker(d) \to 0 \cdots
\]
obtained by replacing ${\sB'}^1=\ker(d)$ by the inverse image, via
$\ker(d)\to \underline\h^1(\sB) = \pic_D(\bar X)$ over $X$, of
$\pic_D^0(\bar X)$ over $X$. Let $\sA'=\sA\cap \sB'$ and let $\sC'$ be the image of
$\sB'$. We have $\sA'=\tau_{\leqslant 1}(\sA)$, and $\sC'$ is deduced from $\sC$
as $\sB'$ from $\sB$, whence a commutative diagram of short exact sequences
\begin{equation*}\tag{3.2.3}\label{V:eq:3-2-3}
\xymatrix{
  0 \ar[r]
    & \sA \ar[r]
    & \sB \ar[r]
    & \sC \ar[r]
    & 0 \\
  0 \ar[r]
    & \tau_{\leqslant 1} \sA \ar[r] \ar[u] \ar[d]
    & \sB' \ar[r] \ar[u] \ar[d]
    & \sC' \ar[r] \ar[u] \ar[d]
    & 0 \\
  0 \ar[r]
    & \h_c^1(X,\dmu_n)[1] \ar[r]
    & \pic_D^0(\bar X)[1] \ar[r]
    & \pic_D^0(\bar X)[1] \ar[r]
    & 0
}
\end{equation*}
where the last row is to be seen as an exact sequence of complexes
of sheaves concentrated in degree $1$ over $X$.
\begin{enumerate}[a)]
  \item $e\in \hh^2(X,\sA)$ comes from
    $\widetilde e\in \hh^2(X,\tau_{\leqslant 1}\sA)$ (since its image in
    $E^{02}$ is zero). We seek the image $\bar e$ of $\widetilde e$ in
    $\hh^2(X,\h_c^1(X,\dmu_n)[-1]) = \h^1(X,\h_c^1(X,\dmu_n))$.
  \item $e_1\in \hh^1(X,\sC)$ comes from $\widetilde e_1\in \hh^1(X,\sC')$,
    since its image $f_0$ in $\h^0(X,\pic_D(\bar X))$ lies in
    $\h^0(X,\pic_D^0(\bar X))$. We may take
    $\widetilde e=\partial \widetilde e_1$.
  \item We hence have $\bar e = \partial f_0$, where $\partial $ is defined by
    \eqref{V:eq:3-2-2}, and where we use the usual isomorphisms
    $\hh^2(X,\sF[-1]) = \h^1(X,\sF)$. These isomorphisms anticommute with
    $\partial$, so that $\bar e=-\partial f_0$, for $\partial$ defined
    by the exact sequence
    \[\xymatrix{
      0 \ar[r]
        & \h_c^1(X,\dmu_n) \ar[r]
        & \pic_D^0(\bar X) \ar[r]
        & \pic_D^0(\bar X) \ar[r]
        & 0
    }\]
    We conclude by (\hyperref[IV]{Cycles}, \ref{IV:1-1-1}).
\end{enumerate}




\subsection{}\label{V:3-3}

For every homomorphism $\varphi:\h_c^1(X,\dmu_n)\to \dZ/n$, let
$\varphi(u)\in \h^1(X,\dZ/n)$ be the image of $u$ under
$\varphi:\h^1(X,\h_c^1(X,\dmu_n)) \to \h^1(X,\dZ/n)$. If
$x\in \h_c^1(X,\dmu_n)$, the cup-product $\varphi(u)\smallsmile x$ is in
$\h_c^2(X,\dmu_n)$. We propose to prove that




\begin{proposition_}\label{V:3-4}
$\tr(\varphi(u)\smallsmile x) = \varphi(x)$.
\end{proposition_}

This identity is deduced, by applying $\varphi$, from
\begin{equation*}\tag{3.4.1}\label{V:eq:3-4-1}
  \tr(u\smallsmile x) = x
\end{equation*}
where $\tr$ is this time the map
\[\xymatrix{
  \tr:\h_c^2(X,\h_c^1(X,\dmu_n)\otimes\dmu_n) \ar[r]
    & \h_c^1(X,\dmu_n) \text{,}
}\]
so that $x\mapsto \tr(u\smallsmile x)$ is the composite map
\small
\begin{equation*}\tag{3.4.2}\label{V:eq:3-4-2}
\begin{aligned}
&\xymatrix{
  \h_c^1(X,\dmu_n) \ar[r]^-{u\otimes x}
    & \h^1(X,\h_c^1(X,\dmu_n)) \otimes \h_c^1(X,\dmu_n)
    & \ar[l]_-\sim \h^1(X,\dZ/n)\otimes \h_c^1(X,\dmu_n) \otimes \h_c^1(X,\dmu_n)} \\ &\xymatrix{
    \ar[r]^-{(1)}
      & \h_c^2(X,\dmu_n) \otimes \h_c^1(X,\dmu_n) \ar[r]^-{\tr\otimes\operatorname{id}}
      & \h_c^1(X,\dmu_n)
}
\end{aligned}
\end{equation*}
\normalsize
where (1) is given by
$a\otimes b\otimes c\mapsto (a\smallsmile c)\otimes b$.

Let us express this morphism in terms of the trace morphism
$\tr:\h_c^3(X\times X,\dmu_n^{\otimes 2}) \to \h_c^2(X,\dmu_n)$ relative to the
second projection. For $x\in \h_c^3$, $\tr(x)$ is computed as follows: among its
K\"unneth components, we keep only that of type $(2,1)$, and we apply
$\tr:\h_c^2(X,\dmu_n) \to \dZ/n$.

Denote by $u'$ the image of $u$ under the map
\begin{align*}
&\xymatrix{
  \h^1(X,\h_c^1(X,\dmu_n))
    & \ar[l]_-\sim \h^1(X,\dZ/n) \otimes \h_c^1(X,\dmu_n) = \h^1(X,\dZ/n)\otimes \h^1(\bar X,j_!\dmu_n)
} \\
&\xymatrix{
  \ar[r]
    & \h^2(X\times \bar X,\dZ/n\boxtimes j_!\dmu_n) \text{.}
}
\end{align*}
In view of the relation between K\"unneth and cup-product, we have
\[
  \tr(u\smallsmile x) = -\tr(u'\smallsmile \pr_1^\ast x) \text{.}
\]
The second cup-product is that of (\hyperref[IV]{Cycle} \ref{IV:1-2-4}):
\[\xymatrix{
  \h^2(X\times \bar X,\dZ/n\boxtimes j_!\dmu_n) \otimes \h^1(\bar X\times \bar X,j_!\dmu_n\boxtimes \dZ/n) \ar[r]
    & \h^3(\bar X\times \bar X,j_!\dmu_n \boxtimes j_!\dmu_n) \text{.}
}\]
The sign -- comes from permuting two symbols of degree $1$ in
(1). Since $c\smallsmile\pr_1^\ast x$ and
$c^{1,1}\smallsmile \pr_1^\ast$ have the same $(2,1)$ K\"unneth
component, and since $u'=c^{1,1}$, we have
\[
  \tr(u\smallsmile x) = \tr(c\smallsmile \pr_1^\ast x) \text{.}
\]
The cup-product on the right may this time be interpreted as the
cup-product (\hyperref[IV]{Cycle} \ref{IV:1-2-5}) of
$\cl(\Delta)\in \h_\Delta^2(X\times X,\dmu_n)$ with
$\pr_1^\ast x\in \h_c^1(\Delta,\dmu_n)$. On $\Delta$, we have
$\pr_1=\pr_2$, whence
\[
  \tr(c\smallsmile \pr_1^\ast x) = \tr(\cl(\Delta)\smallsmile \pr_1^\ast x) = \tr(\cl(\Delta)\smallsmile \pr_2^\ast x) = \tr(\cl(\Delta))\smallsmile x = x \text{.}
\]









\section{Compatibility SGA 4 XVIII 3.1.10.3}\label{V:4}

In \cite[XVIII]{sga4}, in the grip of discouragement, I did not check
the compatibility of the title. Although it turns out to be mere smoke, I feel
obliged to repair this gap here. For $f:X\to S$ separated of finite
type, the point was to check that the adjunction morphism
$\eR f_!\eR f^!\sL\to \sL$ is compatible with every \'etale localization
$k:V\to S$. Denoting by $k$ and $f$ several arrows, as below,
the point is to check a commutativity:
\[\xymatrix{
  X_V \ar[r]^-k \ar[d]^-f
    & X \ar[d]^-f \\
  V \ar[r]^-k
    & S
} \qquad
\xymatrix{
  k^\ast \eR f_! \eR f^!\sL \ar[r]^-{k^\ast(\text{adj})} \ar[d]^-\sim
    & k^\ast \sL \ar@{=}[d] \\
  \eR f_! \eR f^! k^\ast \sL \ar[r]^-{\text{adj}}
    & k^\ast \sL \text{.}
}\]
At the level of complexes, the desired commutativity reads (loc.\ cit.)
\[\xymatrix{
  k^\ast f_!^\bullet f^!_\bullet \sL \ar[r] \ar[d]^-{(1)}
    & k^\ast \sL \ar@{=}[d] \\
  f_!^\bullet f^!_\bullet k^\ast \sL \ar[r]
    & k^\ast \sL \text{.}
}\]
We check it component by component, which reduces to a
compatibility for $\sL$ a sheaf and for each of the pairs of adjoint
functors $(f_i^!,f_!^i)$, simply denoted $f^!$ and $f_!$. The arrow (1)
is defined from the isomorphism $(2):k^\ast f_!=f_! k^\ast$, and from a
morphism $(3):k^\ast f^! \to f^! k^\ast$ deduced from it. In loc.\ cit., one
first introduces a morphism $(4):k_! f_!\to f_! k_!$, deduced from $(2)$ by
adjunction of $k^\ast=k^!$ and of
$k_!:k_! f_!\to k_! f_! k^\ast k_!\to k_! k^\ast f_! k_! \to f_! k_!$, and one
deduces $(3)$ from $(4)$ by adjunction. It amounts to the same to deduce $(3)$
from $(2)$ by adjunction of $f_!$ and
$f^!:k^\ast \to f^! f_! k^\ast f^!\to f^! k^\ast f_! f^! \to f^! k^\ast$.

Writing $k$ for $k^\ast$ and omitting $\sL$, the commutativity is then
that of the outer border of
\[\xymatrix@=.5cm{
  k f_! f^! \ar[r] \ar@{=}[d]
    & f_! k f^! \ar[r]
    & f_! f^! f_! k f^! \ar[r]
    & f_! f^! k f_! f^! \ar@{=}[d] \\
  k f_! f^! \ar[d]
    & & & \ar[lll] f_! f^! k f_! f^! \ar[d] \\
  k
    & & & \ar[lll] f_! f^! k \text{.}
}\]
In the first row, one finds a homomorphism
$k _! f^! \to f_! f^! k f_! f^!$, defined because $k f_! f^!$ is of the
form $f_! X$ ($X=k f^!$); it admits as retraction the
adjunction morphism in the second row, and the commutativity of the lower square is
the functoriality of this adjunction morphism.



